\scanpage{173}

\uline{Le foncteur de Picard}

\margin{Faits élémentaires sur les classes de diviseurs, faisceaux inversibles, leurs propriétés etc}

Soit $X$ un $S$-préschéma, on pose

\[
\mathrm{Prépic}\,(X/S)=\deleted{H^1(X,\underline O_X^*)}\mathrm{Pic}(X\deleted{\unsure{/S}})=H^1(X,\underline O_X^*)
\]

\[
\text{\uuline{Prépic}}_{X/S}(S')=\mathrm{Prépic}\,(X'/S')=H^1(X',\underline O_{X'}^*)\qquad\text{où }\ X'=X\times_S S'
\]

donc $\text{\uuline{Prépic}}_{X/S}$ est bien un foncteur \uuline{contravariant} en le préschéma $S'$

\[
\mathrm{Prépic}'(X/S)=\Gamma(S,\underline H^1_{X/S}(\underline O_X^*))
\]

\[
\text{\uuline{Prépic}}'_{X/S}(S')=\Gamma(S',\underline H^1_{X'/S'}(\underline O_{X'}^*))\qquad\text{où }\ X'=X\times_S S'
\]

% slanted
\margin{expliciter les objets, l'\unsure{équivalence}, \unsure{pour} $\mathrm{Pic}(X\deleted{/S})$ définir \unsure{fn} $\mathrm{Pic}(X\deleted{/S})\longrightarrow\mathrm{Prépic}'(X/S)$}

On passe de \uuline{Prépic} à \uuline{Prépic}$'$ en exigeant la \unsure{descente} par recouvrements ouverts.

\uline{Proposition 1} Supposons $\underline H^0_{X/S}(\underline O_X)=\underline O_S$. Alors on a un \deleted{isomorphisme}\inserted{\uline{monomorphisme}} canonique

% slanted
\margin{mettre \deleted{\unsure{ce} corollaire} sous forme d'une suite exacte plutôt}

\[
\mathrm{Prépic}'(X/S)\mathrel{\deleted{\unsure{\sim}}\!\longleftarrow}\mathrm{Prépic}\,(X/S)/\mathrm{Pic\,abs}\,(S)
\]

qui est un isomorphisme si $X/S$ a une section Résulte de l'assertion plus précise

\uline{Corollaire 1} Le foncteur $\underline F\longrightarrow\underline F\otimes_S X$ des faisceaux loc. libres de $S$ dans les faisceaux loc. libres de $X$ \deleted{définit un}\inserted{est} \uline{pleinement} fidèle, l'image essentielle est formée des faisceaux loc. libres sur $X$ qui sont triviaux loc. au dessus de $Y$, le foncteur inverse sur ce dernier est le foncteur $\underline H^0_{X/\unsure{S|B}}$.

% slanted
\margin{\unsure{cela} pour un \unsure{cert.|cond:} \unsure{groupe|graphe} ? ou une proposition préliminaire}

% slanted
\margin{dire que les \unsure{images} sont \unsure{plus} \unsure{petits gr.}}

% archive sheet number: 172
